\documentclass[11pt]{article}
\usepackage[margin=1in]{geometry}
\usepackage{amsmath,amssymb,amsthm}
\usepackage{booktabs}
\usepackage{hyperref}
\usepackage{array}

\newtheorem{theorem}{Theorem}[section]
\newtheorem{proposition}[theorem]{Proposition}
\newtheorem{lemma}[theorem]{Lemma}
\newtheorem{corollary}[theorem]{Corollary}
\newtheorem{definition}[theorem]{Definition}
\newtheorem{conjecture}[theorem]{Conjecture}
\newtheorem*{openproblem}{Open Problem}

\title{How large is the equality-of-factors automaton?\\
Empirical evidence and construction strategy for an open problem of Khodier}
\author{Andrew Hingston}
\date{\today}

\begin{document}
\maketitle

\begin{abstract}
For a $k$-automatic sequence $T$ produced by an $m$-state DFAO, the \emph{equality of
factors} predicate $FE(i,j,l) := \forall t\,(t<l \Rightarrow T[i+t]=T[j+t])$ defines a
language whose minimal automaton size $|FE|$ is of interest both computationally (Walnut
and related tools routinely blow up while constructing it) and structurally: Khodier's
2026 thesis states as Open Problem 1 that no class of examples with \emph{exponential}
blowup in $m$ is known, and conjectures that one exists. We report a large empirical
sweep (random $k$-uniform morphisms, $k\in\{2,3\}$, $m\le7$) that finds no such family in
the bulk ($|FE|$ scales like $2$--$4\,m^3$ at the median, with a small unresolved tail),
identify that most apparent construction-time blowups are digit-order artefacts rather
than genuine largeness of the minimal automaton, isolate ``lossy coding of a group
automaton'' as the one lead that empirically inflates $|FE|$, and give a
reverse-first-determinization construction strategy that solves several inputs on which a
current build of Walnut fails at equal memory. On the theory side we report an
unconditional, referee-checked result: for the generalised Thue--Morse family $G_p$ (a
$p$-state DFAO), $|FE_{G_p}|_{\mathrm{msd}}=\Theta(p^3)$, exactly $p^3+8$ for
$3\le p\le24$ -- the first infinite family with $|FE|$ pinned to a matching polynomial
order. We also report a conditional upper bound
$|FE|\le m^4+m^6+m^8\Lambda(T)^2$ in terms of a computable, morphism-only invariant
$\Lambda$, and state honestly what remains open: no exponential family, and no
unconditional upper bound smaller than the literature's $2^{9m^2}$, has been found.
\end{abstract}

\tableofcontents

\section{Introduction and problem statement}

Automatic sequences -- infinite sequences produced by finite automata reading the
base-$k$ digits of $n$ -- support a decision procedure (implemented in tools such as
Walnut) that reduces first-order statements about the sequence to automaton
constructions. A recurring primitive in that procedure is the \emph{equality of
factors} predicate
\[
FE(i,j,l) \;:=\; \forall t\,\big(t<l \Rightarrow T[i+t]=T[j+t]\big),
\]
asserting that the length-$l$ factors of $T$ starting at positions $i$ and $j$ agree.
$FE$ underlies statements about repetitions, palindromes, critical exponents, and
recurrence, and its automaton is frequently built as an intermediate step in larger
queries even when it is not the final object of interest.

Mazen Khodier's 2026 PhD thesis (\emph{New Methods for Analyzing the Properties of
Automatic Sequences}, University of Waterloo, Chapter 8) records that the size of the
minimal automaton for $FE$, relative to the size $m$ of the DFAO producing $T$, is open.
Verbatim:

\begin{openproblem}[Khodier 2026, Chapter 8]
Give a characterization for the $k$-automatic sequences that have an exponential blowup
in size when constructing the ``equality of factors'' automaton. Currently no such class
of examples is known.
This issue has two distinct aspects. The first concerns the size of the minimal
automaton for the equality of factors predicate relative to the size of the morphism.
We believe that this relationship is exponential, but currently, no class of examples
is known that proves this behavior. The second aspect involves determining the most
efficient strategy for constructing this automaton [\ldots]
\end{openproblem}

The thesis records one existing upper bound, for a different but logically equivalent
formulation
$$FE_2(i,j,n) := \forall u,v.\ (u\ge i \,\wedge\, u<i+n \,\wedge\, u+j=v+i) \Rightarrow T[u]=T[v],$$
namely $|FE_2|\le 2^{9m^2}$, and a worked anecdote: the Tribonacci word's minimal $FE$
has only $26$ states, but the naive one-shot Walnut construction of it peaked at
$323{,}831{,}403$ intermediate states and $300$~GB before finishing. That anecdote is the
seed of the present study: it shows the two aspects of the open problem are decoupled --
the \emph{final} automaton can be small even when the \emph{construction} is enormous --
and raises the question of how general that decoupling is.

This paper reports on an empirical and partly theoretical investigation of both aspects,
using a purpose-built Rust automaton engine (\S\ref{sec:method}) with a Brzozowski
reverse-first construction strategy and hard memory guards, run at scale over random and
structured families of automatic sequences. \S\ref{sec:method} describes the engine and
sweep design. \S\ref{sec:results} reports the random-ensemble table, the digit-order
artefact finding, structured-family results, and the coding-spectrum lead. \S\ref{sec:construction}
gives the construction strategy and a like-for-like benchmark against Walnut.
\S\ref{sec:theory} states what is unconditionally proved and what remains conjectural,
following an independent adversarial referee pass (\texttt{paper/proof-verdict.md}).
\S\ref{sec:conjecture} states the polynomial conjecture with an honest confidence level.
\S\ref{sec:limitations} and \S\ref{sec:repro} close with limitations and reproducibility.

\emph{On overclaiming.} Every number in this paper traces to a file under
\texttt{results/*.json}, a table in \texttt{docs/TARGET1.md} or \texttt{docs/TARGET1-TABLE.md},
or a lemma/theorem in \texttt{paper/proof-*.md} whose status has been independently
checked (\texttt{paper/proof-verdict.md}). Section~\ref{sec:theory} includes \emph{only}
what the referee verdict marks PROVED; everything the verdict marks fitted, heuristic, or
partially machine-verified is reported as such, not as a theorem.

\section{Method}
\label{sec:method}

\subsection{Engine}
The automaton engine (\texttt{engine/}, Rust) is driven by a small command language over
stdin: \texttt{mode msd|lsd} selects most- or least-significant-digit-first reading;
\texttt{def T k m start w\_0..w\_\{m-1\} coding} defines a $k$-uniform, $m$-state DFAO
with an output coding; \texttt{let NAME(args) formula} defines a named predicate in a
first-order language over $\forall/\exists$, indexing $T[i+t]$, and boolean connectives
\texttt{\& | => !=}, with \texttt{\$NAME(...)} for predicate composition; \texttt{?
formula} decides a closed formula; \texttt{mem} reports live/peak allocator bytes.
$|FE|$ throughout is the number of states of the minimal \emph{complete} DFA, including
the dead state (Walnut's convention differs by exactly $1$; this is corrected for in all
Walnut comparisons in \S\ref{sec:construction}).

\subsection{Guard}
A single earlier sweep run crashed the host machine by launching many unbounded engine
processes in parallel. Three independent guards now prevent this (\texttt{docs/GUARD.md}):
(1) a hard allocator budget inside the engine (\texttt{AM\_MEM\_MB}, default $2048$~MB;
on exceeding it the process exits cleanly with status~3 and no partial state);
(2) \texttt{explore/engine.py}, the sole sanctioned way to launch the engine from Python,
providing admission control (a job does not start until free RAM exceeds a floor and
kernel memory pressure is normal), a watchdog that kills over-budget children, and
guaranteed cleanup on exit; (3) a system-level LaunchAgent that kills any engine process
exceeding an absolute RSS ceiling regardless of how it was launched. A \texttt{FAIL
budget} outcome in a sweep is a datum -- the automaton genuinely exceeded the ceiling --
not an artefact of the harness.

\subsection{Sweep design}
The primary random sweep (\texttt{explore/blowup.py}, plus rescue passes) draws
$k$-uniform admissible morphisms with a binary coding from a sampler named PRISM, for
$k\in\{2,3\}$ and DFAO size $m\in\{2,\ldots,7\}$, targeting up to $40$ admissible
sequences per $(k,m)$ cell. For each sequence, $|FE|$ is measured as the minimal DFA size
of $FE(i,j,l)$. A sequence that fails (memory budget or timeout) in msd is retried in
lsd, and vice versa; sequences that fail in both digit orders are further retried under
Khodier's alternative $FE_2$ formulation and under a forced small-subset-cap ladder
(\S\ref{sec:construction}). The final merged table (\texttt{docs/TARGET1-TABLE.md},
built by \texttt{explore/final\_table.py} from \texttt{results/blowup.json},
\texttt{results/blowup\_retry.json}, \texttt{results/blowup\_residue.json},
\texttt{results/blowup\_residue2.json}, \texttt{results/blowup\_residue3.json}) keeps one
number per sequence, preferring an msd value (msd and lsd sizes for the same predicate
genuinely differ) and falling back to an lsd-only or still-censored classification when
no msd value was obtained in any pass.

Beyond the random ensemble, structured families are swept by dedicated scripts
(\texttt{explore/thin\_families*.py}, \texttt{struct\_families.py}, \texttt{gtm\_full.py},
\texttt{blowup\_features.py}) covering thin/counting sets (fixed number of $1$-bits),
pattern-containment sets, digit-sum-mod-$p$ sequences under both the faithful group
coding and a lossy binary coding, and a feature-correlation study relating prefix
statistics (subword complexity growth, run length, right-special count) to $|FE|$ within
a size class.

\section{Results}
\label{sec:results}

\subsection{Random ensemble: no exponential family in the bulk}

Table~\ref{tab:ensemble} reproduces the final uncensored per-$(k,m)$ quartile table
(\texttt{docs/TARGET1-TABLE.md}), merged over five rescue passes. Of $412$ dispatched
sequences, $356$ resolved to an msd value, $45$ resolved only in lsd (not directly
comparable, hence excluded from the quartiles), and $11$ remain still-censored after five
rescue passes -- all in the two hardest cells, $k=2,m=6$ (one sequence) and $k=3,m=7$
(ten sequences).

\begin{table}[h]
\centering
\caption{Random-ensemble $|FE|$ (msd) by $(k,m)$, quartiles over resolved sequences.
$n$ = dispatched, cens = still-censored. Source: \texttt{docs/TARGET1-TABLE.md},
\texttt{results/final\_table.json}.}
\label{tab:ensemble}
\begin{tabular}{@{}rrrrrrrrrrr@{}}
\toprule
$k$ & $m$ & $n$ & cens & min & q1 & med & q3 & max & med$/m^3$ & $\log_2(\max)/m$ \\
\midrule
2 & 2 & 8  & 0  & 3   & 6   & 8   & 10   & 15   & 0.94 & 1.95 \\
2 & 3 & 33 & 0  & 9   & 34  & 51  & 96   & 218  & 1.89 & 2.59 \\
2 & 4 & 24 & 0  & 24  & 58  & 161 & 404  & 549  & 2.52 & 2.28 \\
2 & 5 & 39 & 0  & 97  & 260 & 340 & 403  & 924  & 2.72 & 1.97 \\
2 & 6 & 40 & 1  & 201 & 474 & 608 & 966  & 2124 & 2.81 & 1.84 \\
2 & 7 & 40 & 0  & 200 & 685 & 882 & 1134 & 1770 & 2.57 & 1.54 \\
3 & 2 & 28 & 0  & 3   & 7   & 17  & 22   & 38   & 2.12 & 2.62 \\
3 & 3 & 40 & 0  & 10  & 50  & 72  & 100  & 190  & 2.69 & 2.52 \\
3 & 4 & 40 & 0  & 135 & 216 & 250 & 348  & 532  & 3.91 & 2.26 \\
3 & 5 & 40 & 0  & 112 & 328 & 424 & 509  & 1604 & 3.40 & 2.13 \\
3 & 6 & 40 & 0  & 311 & 540 & 752 & 872  & 2356 & 3.48 & 1.87 \\
3 & 7 & 40 & 10 & 521 & 722 & 847 & 1085 & 3480 & 2.47 & 1.68 \\
\bottomrule
\end{tabular}
\end{table}

Power-law fits to the medians beat exponential fits by roughly an order of magnitude in
log-residual sum of squares (\texttt{docs/TARGET1-TABLE.md}):
\begin{align*}
k=2:&\quad \mathrm{med} \approx 0.666\, m^{3.81}\ (\mathrm{SSE}=0.18) \qquad\text{vs.}\qquad
        \mathrm{med} \approx 2.44 \cdot 2.496^{m}\ (\mathrm{SSE}=1.26),\\
k=3:&\quad \mathrm{med} \approx 2.11\, m^{3.23}\ (\mathrm{SSE}=0.21) \qquad\text{vs.}\qquad
        \mathrm{med} \approx 6.42 \cdot 2.168^{m}\ (\mathrm{SSE}=1.07).
\end{align*}
$\log_2(\max)/m$, the quantity that would grow without bound under a true exponential
family, instead \emph{falls} with $m$ in five of six columns per $k$. The largest msd
value found anywhere in the ensemble is $|FE|=3480$ ($k=3,m=7$,
\texttt{def T 3 7 0 044 513 202 604 520 421 061 1111110}), surfaced only via the forced
small-cap ladder of \S\ref{sec:construction} -- first-pass msd and every other
formulation failed on it.

This is evidence about the \emph{typical} case under one random-morphism sampler, not a
proof about all automatic sequences; a single adversarial family would settle Open
Problem 1(A) regardless of what typical sequences look like. The $11$ still-censored
sequences, concentrated entirely in $k=3,m=7$, are exactly where such a family would have
to live and remain unresolved.

\subsection{Finding: most apparent blowups are digit-order artefacts}

Of $47$ sequences that failed the memory budget in msd during the first sweep pass,
$45$ (96\%) finished in lsd at modest size ($64$--$989$ states for $k=3$, $m\in\{3,4\}$).
Every $k=3$, $m\in\{3,4\}$ failure was such an artefact: the \emph{construction} blew up
in one digit order while the \emph{minimal} automaton was small and reachable in the
other. This reproduces, at scale across hundreds of sequences, exactly the pattern of
Khodier's Tribonacci anecdote (26 final states, 323 million intermediate states) --
suggesting the decoupling between construction cost and final size is common, not a
one-off. It also means aspect (B) of Open Problem 1 (construction strategy) is a real and
separate problem from aspect (A) (minimal size): most of what looks like evidence for
exponential blowup in a naive single-pass sweep is not evidence for aspect (A) at all.

\subsection{Structured families}

Table~\ref{tab:struct} summarizes the structured-family sweeps
(\texttt{docs/TARGET1.md}, ``Structured families''). All values are minimal $|FE|$, msd
unless stated; a dash marks an engine failure at the stated budget.

\begin{table}[h]
\centering
\caption{Structured families, $|FE|$ vs.\ parameter. Source: \texttt{docs/TARGET1.md}.}
\label{tab:struct}
\begin{tabular}{@{}lll@{}}
\toprule
Family & Mode & Growth \\
\midrule
exactly $c$ ones, base 2 ($m=c+2$) & msd & 52, 183, --, 1042, 2008, 3463 ($\sim c^3$) \\
exactly $c$ ones, base 2 & lsd & 41, 242, 1133, 4227, 13461, 38504 (ratios $5.9,4.7,3.7,3.2,2.9$) \\
at most $c$ ones, base 2 & msd & 51, 108, --, 327, 502, 730 ($\sim c^{2.5}$) \\
contains $1^r$ in binary ($m=r+1$) & msd & 7, 91, --, --, 998, 1713, 2700, 4004 (ratios $\sim1.5$) \\
contains $10^{r-1}$ & msd & 13, 54, --, --, 587, 1026, 1651 (polynomial) \\
$s_2(n)\bmod p$, full output ($m=p$) & msd & 15, 35, --, 133, 224, 351, 520, 737 ($\sim p^{2.5}$) \\
$s_2(n)\bmod p$, full output & lsd & 22, 39, 56, 77, 102, 131, 164, 201 (exactly quadratic) \\
$[s_2(n)\bmod p \ne 0]$, binary code & msd & 15, 190, 698, 1877, --, --, --, -- (ratios $12.7,3.7,2.7$) \\
\bottomrule
\end{tabular}
\end{table}

Reading (\texttt{docs/TARGET1.md}): thin/counting sets and pattern-containment sets both
give polynomial growth. The generalised-Thue--Morse family with \emph{full} output is
polynomial and exactly quadratic in lsd; the naive heuristic that $FE$ forces
$i\equiv j\pmod{2^p}$ (hence $2^p$ msd states) is refuted by direct measurement and
retracted in that document -- the congruence is forced globally along the whole window,
so the msd automaton never needs a sliding window. The one family that grows fast within
the range measured is $[s_2(n)\bmod p\ne0]$, the \textbf{lossy binary coding} of the same
$p$-state group automaton, discussed next.

\subsection{The coding-spectrum lead}

Collapsing $s_2(n)\bmod p$ from a faithful $p$-letter output to the single bit
$[s_2(n)\not\equiv0\bmod p]$ keeps the DFAO at $p$ states but multiplies $|FE|$ by
roughly $100\times$ at $p=4$ in lsd ($6154$ vs.\ $56$; \texttt{docs/TARGET1.md}). This is
the one effect in the entire sweep where a size-preserving change to the sequence
(recoding the output alphabet, not changing the transition structure) produces a large
multiplicative jump in $|FE|$, and it is the empirical seed of the theoretical work in
\S\ref{sec:theory}: the relevant invariant is not ``big automaton'' but ``lossy coding of
a group automaton,'' where the $FE$ automaton must track a pair of group elements rather
than their difference.

\subsection{The $3p^4$ fit, and its retraction}

An earlier fit to the singleton-coding family $T_p=[s_2(n)\equiv1\bmod p]$ reported
$|FE_{T_p}|\approx 3p^4$ to within $2.5\%$ using six points, $p=3,\ldots,8$: $190, 698,
1877, 3971, 7243, 11988$. This fit is \textbf{not} carried forward as established. The
referee review (\texttt{paper/proof-verdict.md}, on \texttt{proof-lower.md} \S5) and the
family document itself (\texttt{proof-family.md} \S5.4) both flag it as unsafe: at
$p=3$, $3p^4=243$ against a measured $190$, a $28\%$ error, not $2.5\%$; and the local
log-log slope of the same six points \emph{falls through} $4$ (from $4.52$ to $3.77$
over $p=3..8$), while a cubic fit ($44.4p^3-205.4p^2+308.5p-88.7$) matches the same data
with residual $\le5.7$. The decisive point that would separate a cubic from a quartic
law, $p=9$ msd, does not fit in a $6$~GB memory guard and remains unresolved (confirmed
again during the referee pass: allocator budget exceeded, return code~3). The status of
$T_p$'s growth exponent is therefore \textbf{open}, not $\Theta(p^4)$ as earlier stated.

\section{Construction strategy and the Walnut benchmark}
\label{sec:construction}

\subsection{The ladder}
Across every hard case encountered, a small forward-subset-construction cap (e.g.\
$50{,}000$ subsets) that fails fast into a Brzozowski reverse-first determinization
consistently beats a large forward cap by orders of magnitude: an exact-$c=3$ counting
sequence resolves to $482$ states in effectively no time under the small-cap-then-reverse
ladder, versus exhausting a $6$~GB budget under a large forward cap; $[s_2\equiv a
\bmod 5]$ resolves at $30$~MB versus blowing $3.3$~GB (\texttt{docs/TARGET1.md}). This
ladder -- small forward cap, Brzozowski reverse, larger forward cap, Brzozowski reverse
again -- is now the engine's default construction strategy.

\subsection{Like-for-like benchmark against Walnut}
\texttt{bench/README.md} reports a controlled comparison against Walnut~8-dev (head
2026-06-15, OpenJDK~26, \texttt{java -Xmx6g}), both run with a 6~GB memory ceiling and,
for the engine, a 15-minute wall-clock timeout on the Walnut side. Both tools count
minimal DFA states; the engine's count includes the dead state, Walnut's does not, so
\texttt{ours = Walnut + 1} whenever both finish -- itself an independent cross-check of
every $|FE|$ number reported elsewhere in this paper.

\begin{table}[h]
\centering
\caption{$|FE|$ and wall time, engine vs.\ Walnut 8-dev, 6~GB each. Source:
\texttt{bench/README.md}.}
\label{tab:walnut}
\begin{tabular}{@{}lrrrr@{}}
\toprule
Sequence & Ours (states) & $s$ & Walnut (states) & $s$ \\
\midrule
Thue--Morse & 15 & 0.0 & 14 & 0.3 \\
period-doubling & 8 & 0.0 & 7 & 0.2 \\
Rudin--Shapiro & 68 & 0.1 & 67 & 0.4 \\
paperfolding & 44 & 0.0 & 43 & 0.2 \\
Cantor & 17 & 0.0 & 16 & 0.2 \\
Mephisto & 14 & 0.0 & 13 & 0.3 \\
PRISM-a & 24 & 0.0 & 23 & 0.3 \\
PRISM-d & 82 & 0.5 & 81 & 0.9 \\
champion-$m5$ & 199 & 0.0 & 198 & 6.1 \\
$k3m3$-artefact-b & 71 & 0.0 & 70 & 18.0 \\
$k3m3$-artefact-a & 216 & 0.0 & 215 & 444.5 \\
PRISM-1 ($k{=}4,m{=}6$) & 467 & 38.6 & OOM & 251.5 \\
$[s_2\equiv a\bmod3]$ & 190 & 1.5 & timeout & 900 \\
$[s_2\equiv a\bmod4]$ & 698 & 0.2 & OOM & 313 \\
$[s_2\equiv a\bmod5]$ & 1877 & 2.1 & OOM & 92 \\
$[s_2\equiv a\bmod6]$ & 3971 & 20.9 & OOM & 53 \\
tail-a ($k{=}2,m{=}7$) & 1165 & 139 & OOM & 77 \\
tail-b ($k{=}3,m{=}5$) & 1000 & 168 & OOM & 119 \\
tail-c ($k{=}2,m{=}6$) & fail (535) & -- & OOM & 164 \\
\bottomrule
\end{tabular}
\end{table}

On easy inputs the two tools are equivalent -- both instant, same automaton. On the
FE-hard inputs (group automata with lossy codings, PRISM-1, the sweep tail), Walnut
8-dev at $6$~GB fails on $7$ of $9$ and needs $7$ minutes on an eighth, while the engine
answers $8$ of $9$ in $0$--$170$~s at the same memory. The difference, per
\texttt{bench/README.md}, is construction strategy -- reverse-first determinization with
a small forward cap, and the lsd/msd switch -- not raw speed; Walnut has a
\texttt{reverse} command and could adopt the same strategy, but it is not its default.

\paragraph{Fairness caveats, quoted verbatim from \texttt{bench/README.md}.} ``Walnut got
its default forward msd construction with no expert reformulation (Khodier's
self-verifying predicates would rescue some of these); a bigger heap would rescue some
more; and this is Walnut on the predicate its own authors document as the pathological
one. It is a benchmark of one predicate family, not of the tools.''

\subsection{\texttt{learnfe}: guess-and-verify construction}
\label{sec:learnfe}

Table~\ref{tab:walnut}'s one remaining direct-construction failure, tail-c ($k{=}2,m{=}6$,
\texttt{def T 2 6 0 05 23 44 42 51 10 000010}), which stalled at $535$ states under a
$6$~GB budget, is resolved by a second construction strategy implemented as the engine
command \texttt{learnfe} (\texttt{docs/LEARNFE.md}). Rather than compiling the
universally-quantified formula $FE(i,j,l)=\forall t\,(t<l\Rightarrow T[i+t]=T[j+t])$
directly -- the projection over $t$ that causes the blowup -- \texttt{learnfe} builds a
candidate DFA by Kearns--Vazirani active learning against a membership oracle
$FE(i,j,l)\Leftrightarrow l\le\mathrm{LCP}(i,j)$, then \emph{verifies} the candidate
against a recurrence
$$H(i,j,l) \iff \big(l=0 \ \text{or}\ (T[i]{=}T[j]\ \text{and}\ H(i{+}1,j{+}1,l{-}1))\big)$$
that only $FE$ satisfies (a one-line induction on $l$; \texttt{docs/LEARNFE.md} \S2), so
the reported automaton is exactly $FE$ regardless of how the candidate was reached --
there is no notion of an approximately-correct answer. On tail-c itself, \texttt{learnfe}
completes at $\mathbf{1382}$ states in $14.2$~s and $50$~MB peak. On the $17$ cases where
both constructions finish, they agree exactly, checked inside the engine as
\texttt{A i,j,l. \$FE(i,j,l) <=> \$G(i,j,l)} $\to$ \texttt{TRUE}
(\texttt{docs/LEARNFE.md} \S6.1), and independently, without any automaton machinery, by a
pure-Python brute force over morphism-generated prefixes diffing every $FE(i,j,l)$,
$i,j,l<B$ against the learned automaton ($416/416$ for Rudin--Shapiro at $B{=}12$ in both
digit orders, $826/826$ for tail-c at $B{=}14$, $325/325$ for tail-b at $B{=}11$).

Because tail-c now has a verified value, the Walnut comparison of
Table~\ref{tab:walnut} becomes $\mathbf{9/9}$: on all nine FE-hard inputs (the four
group-automaton lossy codings, PRISM-1, and the four sweep-tail sequences), the engine
now produces a verified answer at $6$~GB where Walnut 8-dev fails on seven and needs
seven minutes on an eighth.

\texttt{learnfe} also resolves three further sequences from the still-censored tail of
\S\ref{sec:results} that neither digit order of the direct construction could build at a
$2.5$~GB budget in $360$~s (\texttt{results/learnfe\_censored.json}): tail-b
($k{=}3,m{=}5$) at $1000$ states, $25.7$~s, $179$~MB; \texttt{c-3.7d} ($k{=}3,m{=}7$) at
$1147$ states, $44.4$~s, $109$~MB; and \texttt{c-3.7h} ($k{=}3,m{=}7$) at $781$ states,
$23.2$~s, $80$~MB. Nine of the thirteen sequences attempted in that pass remain
unresolved at this budget -- all $k=3$, where the $27$-letter track alphabet makes each
equivalence query an order of magnitude more expensive. \texttt{learnfe} is not
uniformly faster than the direct construction: on easy inputs (e.g.\ Thue--Morse) the
direct construction wins by orders of magnitude, since \texttt{learnfe}'s cost tracks the
\emph{answer} size rather than any intermediate blowup. It is a second tool for aspect
(B) of Open Problem 1, not a replacement for the ladder of \S\ref{sec:construction}.

\section{Towards a proof}
\label{sec:theory}

This section reports only what an independent adversarial referee pass
(\texttt{paper/proof-verdict.md}, checking every load-bearing lemma by hand and, where
finite, by an independent brute force not written by the original authors) marked
\textbf{PROVED}. Results marked as fits, heuristics, or partially machine-verified are
labelled accordingly; none are stated as theorems.

\subsection{Proved: a matching-order family, $|FE_{G_p}| = \Theta(p^3)$}

Let $G_p[n] := s(n) \bmod p$ (the digit sum of $n$ in base $2$, reduced mod $p$), the
fixed point of the $2$-uniform morphism $a\mapsto a\,(a{+}1)$ on $\mathbb Z_p$ with the
identity coding: a $p$-state DFAO ($m=p$), the \emph{generalised Thue--Morse word}.

\begin{theorem}[Referee-verified; \texttt{proof-family.md} Thm.\ 4.4, 4.5]
\label{thm:gp}
For every $p\ge3$, the minimal msd DFA of $FE_{G_p}$ has at least $p^3$ states, and
$$p^3 \;\le\; |FE_{G_p}|_{\mathrm{msd}} \;\le\; 34p^3+98p^2+41p \;=\; O(m^3),\qquad
|FE_{G_p}|_{\mathrm{lsd}} \;\le\; 128p^2+160p+19 \;=\; O(m^2).$$
\end{theorem}
The lower bound rests on a parity-rigidity lemma (a factor of $G_p$ of length $\ge4$
determines the parity of its occurrence positions) and a Myhill--Nerode argument
distinguishing $p^3$ residuals by an explicit construction of prefixes; the upper bound
is by inspection of an explicit finite state space of size $O(p^3)$ (resp.\ $O(p^2)$ in
lsd). Both directions were independently re-derived by hand and, where finite, checked by
brute force ($p=2..7$, up to $2{,}058{,}000$ triples for the characterisation lemma) by
the referee pass, with $0$ mismatches.

Separately, and only for each individual $p$ in the range checked (a machine-verified
statement, not a proof for all $p$):
$$|FE_{G_p}|_{\mathrm{msd}} = p^3+8,\qquad |FE_{G_p}|_{\mathrm{lsd}} = 2p^2+3p+12,\qquad
3\le p\le24$$
(exception $p=2$, ordinary Thue--Morse: $15$ and $22$). Engine values reproduced by the
referee at $p=3,7,11,13,16$: msd $=35,351,1339,2205,4104$; lsd $=39,131,287,389,572$,
matching the closed forms exactly.

This is, to our knowledge, the first infinite DFAO family for which $|FE|$ is pinned to a
matching polynomial order -- evidence \emph{against} exponential blowup for this specific
class, against a literature ceiling of $2^{9m^2}$ (Khodier 2026) for the general case.

\subsection{Proved (conditional): a computable predictor $\Lambda$, and $|FE|\le m^4+m^6+m^8\Lambda(T)^2$}

\begin{theorem}[Referee-verified; \texttt{proof-upper.md} Thm.\ 5.4, and supporting Lemma 3.1, Cor.\ 3.2, Thm.\ 4.3]
\label{thm:lambda}
For any $k$-automatic $T$ with minimal msd DFAO of size $m$, the Myhill--Nerode residual
of an $FE$ prefix $(I,J,L)$ with $L\ge2$ is determined by exactly $8$ DFAO states plus
one further datum $\Theta$ (the \emph{middle language}), which depends on $(I,J,L)$ only
through two subsets of $Q^3$. Consequently there is a morphism-only invariant
$\Lambda(T)$, computable in time polynomial in $m^3$ and $\Lambda$ without ever
constructing $FE$, such that
$$|FE_{\mathrm{msd}}(T)| \;\le\; m^4+m^6+m^8\,\Lambda(T)^2.$$
\end{theorem}
This was checked by \texttt{paper/proof-upper-check.py} with $0$ violations on the
authors' worked examples, and independently re-run by the referee on five adversarial
DFAOs the authors had not tested (an $m=4$ extremal automaton, a thin set, a ruler
family, and both codings of $\mathbb Z_4$), again $0$ violations. Table~\ref{tab:lambda}
gives measured $(\Lambda, |FE|)$ pairs.

\begin{table}[h]
\centering
\caption{$\Lambda$ vs.\ measured $|FE|$ on adversarial examples. Source:
\texttt{paper/proof-verdict.md} \S1.}
\label{tab:lambda}
\begin{tabular}{@{}lrrr@{}}
\toprule
Example & $\Lambda$ & $\#\Theta$ & $|FE|$ (engine) \\
\midrule
$C_4$ champion (\texttt{0 01 22 33 10 / 1110}) & 74 & 61 & 1152 \\
ruler $r=4$ & 35 & 61 & 113 \\
exactly-2-ones & 30 & 10 & 183 \\
$\mathbb Z_4$ singleton coding & 28 & 24 & 698 \\
$\mathbb Z_4$ faithful coding & 4 & 2 & 72 \\
\bottomrule
\end{tabular}
\end{table}

\textbf{What this theorem does and does not give.} $\Lambda$ separates the extremes
($4$ for the faithful group automaton, $74$ for the densest $m=4$ example found), and it
is empirically the only invariant in this study that separates a faithful coding from a
lossy coding of the \emph{same} underlying automaton. But $\Lambda$ is \emph{not}
monotone in $|FE|$ across families (the ruler has $\Lambda=35$, $|FE|=113$; the singleton
$\mathbb Z_4$ coding has $\Lambda=28$, $|FE|=698$ -- larger $|FE|$, smaller $\Lambda$),
and nothing in the theorem bounds $\Lambda$ itself. Since $\Lambda\le2^{\min(\gamma,m^3)}$
for a related quantity $\gamma$ that is not otherwise controlled, the unconditional
content of Theorem~\ref{thm:lambda} is $|FE| = 2^{O(m^3)}$ -- \emph{weaker} than the
literature's $2^{9m^2}$ (Khodier 2026), not an improvement on it. The referee pass is
explicit that Theorem~5.4, taken as a bound, is a reparametrisation of the obstruction,
not a solution; its real value is $\Lambda$ as a cheap, engine-free predictor computed
without ever building $FE$.

\subsection{Not proved: the coding-lead is not closed}

An earlier draft of \texttt{proof-family.md} claimed that a level decomposition of the
lossy-coding mechanism (Prop.\ 5.3: at bit-level $\kappa=v_2(j-i)$ the effective coding is
$\chi_{\min(\kappa,p-1)}$, saturating at the identity coding by level $p-2$) ``closes''
the lossy-coding lead as a route to exponential blowup. The referee pass identifies this
as \textbf{circular}: level $0$ of the grading \emph{is} the original singleton-coded
predicate $FE_{T_p}$ restricted to odd differences, so the argument's conclusion --
``$|FE_{T_p}|$ is not exponential unless some single level is already exponential'' --
reduces to ``$|FE_{T_p}|$ is not exponential unless $|FE_{T_p}|$ is exponential.'' There
is also an unaddressed possibility that the $p$ levels combine multiplicatively rather
than additively (the msd automaton cannot in general read off $\kappa=v_2(j-i)$ from a
bounded prefix), which is exactly the shape an exponential family would take; this does
not happen in the data measured so far, but that is a measurement, not a proof. No upper
bound of any kind for $|FE_{T_p}|$ exists in any of the three source documents. The
correct summary, per the referee verdict, is: the level decomposition is a correct
structural fact about $FE_{G_p}$'s coding, and it is \emph{not} a proof that lossy-coded
group automata cannot be an exponential family.

\subsection{Gaps, stated explicitly}

\begin{itemize}
\item \textbf{Open Problem 1(A) is not answered by this paper.} No family with
provably exponential $|FE|$ relative to $m$ is exhibited, in either direction. Theorem~\ref{thm:gp}
is evidence \emph{against} exponential blowup for one natural family; it is not a
resolution of the general question.
\item \textbf{No upper bound smaller than $2^{9m^2}$ is proved.} Theorem~\ref{thm:lambda}'s
unconditional content is $2^{O(m^3)}$, strictly weaker than the literature bound it was
hoped to improve on.
\item \textbf{$|FE_{T_p}|$'s growth exponent is open}, not $4$ as an earlier fit
suggested (\S\ref{sec:results}); the decisive data point ($p=9$, msd) exceeds the $6$~GB
memory guard.
\item \textbf{The closed forms $p^3+8$ and $2p^2+3p+12$ are machine-verified for
$3\le p\le24$, not proved for all $p$.} The construction underlying them is proved
correct as a design from the exact characterisation (Cor.\ 3.5), but was brute-force
cross-checked only to $p\le6$; values for $17\le p\le24$ rest on a single implementation.
\item \textbf{The extremal-family claim in the random-tail direction is data-limited.}
A separately measured extremal ladder-construction champion at $m=5$
(\texttt{def T 2 5 0 01 22 33 44 10 11110}) resolves at $|FE|=3846$ msd, beating the
$1200$-sequence sample champion ($2415$) by $1.6\times$ and continuing the family beyond
$m=4$ ($1152$); but only three points ($15$, $1152$, $3846$ at $m=3,4,5$) are available,
their consecutive $\log_2$ increments are $4.14, 2.13, 1.74$ -- falling, but slowly enough
that this is not decisive evidence either for or against a bounded-exponent law, and
$m=6,7$ for this specific family remain unmeasured.
\end{itemize}

\section{Conjecture}
\label{sec:conjecture}

\begin{conjecture}
For every $k$-automatic sequence $T$ produced by a minimal $m$-state DFAO,
$|FE(T)|_{\mathrm{msd}} = O(m^4)$.
\end{conjecture}

\textbf{Honest confidence: low-to-moderate, and explicitly not the position of Khodier's
thesis, which states the opposite belief.} In favour: every random-ensemble median
through $m=7$ (\S\ref{sec:results}) is polynomial with exponent $3.2$--$3.8$; every
structured family measured (thin sets, pattern-containment, faithful group codings) is
polynomial; the one proved infinite family (Theorem~\ref{thm:gp}) is exactly cubic; and
$\log_2(\max)/m$ falls, not rises, with $m$ across the whole random ensemble. Against: the
sweep only reaches $m\le7$ and $11$ sequences remain permanently unresolved in exactly the
cell ($k=3,m=7$) where an exponential family would have to live; the one family whose
exponent is empirically \emph{not} settled ($T_p$, singleton coding of a group automaton)
is precisely the family the field's one lead points at, and its exponent could be $3$,
$4$, or -- unproved to be excluded -- unbounded; and no unconditional polynomial upper
bound exists for any family, only the $2^{O(m^3)}$ bound of Theorem~\ref{thm:lambda},
which is weaker than the known $2^{9m^2}$. This conjecture should be read as ``no evidence
found for the exponential belief, in a large but bounded search,'' not as a claim that the
belief is wrong.

\section{Limitations}
\label{sec:limitations}

\begin{itemize}
\item The random ensemble is drawn from one sampler (PRISM) of admissible morphisms and
reaches only $m\le7$; a single adversarial family well outside this distribution would
settle Open Problem 1(A) regardless of the ensemble's shape.
\item $11$ sequences (all in $k=3,m=7$) remain still-censored after five rescue passes and
several hours of machine time; they are the natural place to look next.
\item The construction-strategy benchmark (\S\ref{sec:construction}) compares one
predicate family (equality of factors) between two tools at one memory ceiling; it is
explicitly stated in \texttt{bench/README.md} not to be a general claim of superiority
over Walnut, which received no expert reformulation.
\item The $3p^4$ fit for $T_p$ was retracted after independent review found a $28\%$
error at $p=3$ and a falling local exponent; readers of earlier internal documents
(\texttt{docs/TARGET1.md}) should treat that number as superseded by
\S\ref{sec:results} of this paper.
\item The $\Lambda$ invariant (Theorem~\ref{thm:lambda}) is not itself known to be
polynomial in $m$; its unconditional bound is weaker than the pre-existing literature
bound, and its value here is as a cheap empirical predictor, not a proof technique for
Open Problem 1(A) in either direction.
\end{itemize}

\section{Reproducibility}
\label{sec:repro}

All numbers in this paper trace to committed artefacts.

\begin{itemize}
\item \textbf{Engine:} \texttt{engine/} (Rust; \texttt{cargo build --release}), driven
only via \texttt{explore/engine.py} (\texttt{engine.run} / \texttt{engine.pool}), never
raw \texttt{subprocess} with multiple workers (\texttt{docs/GUARD.md}).
\item \textbf{Random ensemble:} \texttt{explore/blowup.py} (first pass),
\texttt{explore/blowup\_retry.py} (lsd rescue), residue passes for the double-failure
tail; merged by \texttt{explore/final\_table.py} into
\texttt{results/final\_table.json}, tabulated in \texttt{docs/TARGET1-TABLE.md}.
\item \textbf{Structured families:} \texttt{explore/thin\_families.py},
\texttt{explore/thin\_families2.py}, \texttt{explore/struct\_families.py},
\texttt{explore/gtm\_full.py}, \texttt{explore/blowup\_features.py}.
\item \textbf{Walnut benchmark:} \texttt{bench/walnut\_fe.py}, results and fairness
caveats in \texttt{bench/README.md}.
\item \textbf{Theory:} \texttt{paper/proof-upper.md}, \texttt{paper/proof-family.md},
\texttt{paper/proof-lower.md}; machine checks in
\texttt{paper/proof-upper-check.py} and \texttt{paper/verdict-checks/}
(\texttt{fam\_check.py}, \texttt{thm44.py}, \texttt{prop51.py}, \texttt{exh\_nonbin.py});
the full adversarial referee pass is \texttt{paper/proof-verdict.md}, whose one-line
verdict per document and per-lemma status table this section summarizes.
\item \textbf{Guard infrastructure:} \texttt{docs/GUARD.md}, \texttt{explore/memguard.sh}.
\end{itemize}

To reproduce the referee's independent checks:
\begin{verbatim}
cd /Users/andrew/maths
.venv/bin/python paper/proof-upper-check.py
.venv/bin/python paper/verdict-checks/fam_check.py
.venv/bin/python paper/verdict-checks/thm44.py
.venv/bin/python paper/verdict-checks/prop51.py
.venv/bin/python paper/verdict-checks/exh_nonbin.py 3
.venv/bin/python paper/verdict-checks/exh_nonbin.py 4
\end{verbatim}

\section*{Acknowledgements}
This work is a direct response to Open Problem 1 of Mazen Khodier's 2026 thesis, and the
Tribonacci construction anecdote reported there motivated the digit-order-artefact
investigation in \S\ref{sec:results}.

\begin{thebibliography}{9}

\bibitem{khodier2026}
M.\ Khodier, \emph{New Methods for Analyzing the Properties of Automatic Sequences},
PhD thesis, University of Waterloo, January 2026.

\bibitem{kss2025}
M.\ Khodier, L.\ Schaeffer, J.\ Shallit, ``Self-verifying predicates in B\"uchi
arithmetic,'' arXiv:2507.19717 [cs.FL], 2025; also in \emph{IWAA 2026}, LNCS 15981,
pp.\ 237--251.

\bibitem{crs2012}
\'E.\ Charlier, N.\ Rampersad, J.\ Shallit, ``Enumeration and decidable properties of
automatic sequences,'' \emph{International Journal of Foundations of Computer Science}
23 (2012), 1035--1066. arXiv:1102.3698.

\bibitem{schaeffer2013}
L.\ Schaeffer, \emph{Deciding properties of automatic sequences}, MMath thesis,
University of Waterloo, 2013.

\bibitem{gss2013}
D.\ Go\v{c}, L.\ Schaeffer, J.\ Shallit, ``The subword complexity of $k$-automatic
sequences is $k$-synchronized,'' DLT 2013. arXiv:1206.5352.

\bibitem{shallit2023}
J.\ Shallit, \emph{The Logical Approach to Automatic Sequences: Exploring
Combinatorics on Words with Walnut}, LMS Lecture Note Series 482, Cambridge University
Press, 2023.

\end{thebibliography}

\end{document}
